\section{Tests with explicit scientific decisions}
\label{sec:tests}

\subsection{Common measurements and acceptance rules}
For a reference $u$ and prediction $v$, measure raw complex errors
\[
E_s(\xi)=\frac{\norm{v(\xi)-u(\xi)}_{H^s}}
                   {\max(\norm{u(\xi)}_{H^s},\eta_s)},\qquad
\norm{u}_{H^s}^2=L\sum_m(1+k_m^2)^s|\widehat u_m|^2,
\]
where $\eta_s$ is a fixed, declared absolute scale for nearly zero
references. Report $s=0,1,2,3$, maximum amplitude, spectral phase where
the signal is above its floor, and spectra. Phase/translation-aligned
errors are secondary: raw errors retain mistakes in propagation speed
and phase. Record action and modified momentum defects but do not
rank methods by conservation alone.

For a continuous reconstruction $v$, evaluate
$r=\partial_\xi v-F_\epsilon(v)$ on independent propagation stations
and a refined profile grid, including nonlinear modes outside the
output representation. A discrete endpoint defect is labeled separately.
Compare residual norms after the same physical normalization and
reconstruction; a smaller residual need not imply a smaller field error
without a stability estimate.

Proposed first gates are maximum $E_0\leq1\%$, maximum $E_1\leq5\%$,
and errors at most $5\%$ in the prespecified primary event amplitude and
distance. For an event near zero distance, use a declared absolute
distance tolerance rather than an ill-conditioned percentage.
A regularity interpretation also requires $E_2,E_3\leq10\%$ and
derivative resolution. Require reference/refinement and representation
errors to be at least ten times smaller than the corresponding
surrogate tolerance. These are engineering targets to freeze after
the pilot, not universal physical or rigorous error bounds.

Track the upper fifth of the retained Fourier band in each $H^s$ sum.
An initial reference screening target is a weighted-tail fraction below
$10^{-6}$ together with grid agreement; a small tail alone does not pass
a case. If a measured Dysthe--NLS difference is smaller than the
combined numerical and learning error, the surrogate has not resolved
that physical correction. Require those errors below one third of the
effect being interpreted, or report the effect as unresolved.

Define the \emph{qualified prediction distance} as the largest tested
$\Xi$ for which all applicable gates pass at all checked stations up to
$\Xi$. It is a sampled, tolerance-dependent quantity. Distinguish
failure of the surrogate, loss of reference resolution, periodic
wrap-around, and departure from physical-model screens. Increase output
density around peaks and growth; do not infer a continuous supremum
from sparse snapshots.

Report at least three training seeds, individual cases, median and
90th-percentile errors, worst cases and failure counts. Use group-level
uncertainty estimates; correlated stations cannot be counted as
independent observations. A small pilot needs case-level reporting
rather than precise-looking population confidence claims.

\subsection{T0. Verify that every method implements the same equation}
\testgoal{Eliminate sign, phase, nonlocal-operator and data-interface
errors before interpreting a training experiment.}
\testdesign{Check the linear factor $e^{ik^2s}$ and the exact nonlinear
plane wave
\[
u(\xi,\tau)=a\exp\{i[k\tau+(k^2-a^2-6\epsilon a^2k)\xi]\}.
\]
Use a nonconstant two/three-mode field to exercise mean-flow coupling.
Check phase/translation equivariance, Fourier-origin reconstruction,
normalization inversion, and agreement of the NumPy and differentiable RHS.
Check network gradients against directional finite
differences. Reject wrong shapes, duplicated endpoints, invalid hashes,
split leakage and inconsistent actual output intervals.}
\testdecision{Algebraic controls should agree to precision-appropriate
tolerances before training; double-precision low-mode RHS differences
should be near roundoff, with finite-difference gradient errors showing
the expected step-size behavior. A failed control blocks the learning
comparison. The supplied independent checks are only part of T0:
the future differentiable adapter still needs to pass against the core.}

\subsection{T1. Qualify the numerical reference and representation}
\testgoal{Establish that the target trajectories and learning grid
resolve the quantities the research intends to predict.}
\testdesign{Use representative low/high amplitudes, narrow/wide packets,
strong chirps and unstable sidebands. Compare $N,2N,4N$ at fixed $L$,
and $h,h/2,h/4$ on a sufficiently fine grid. Check expected RK4 step
convergence before roundoff dominates. For localized packets compare
$L,2L$ and, if necessary, $4L$ at fixed spacing, reconstructing the same
localized field instead of periodically repeating it. Compare common
windows, derivative tails and mean flow. Perform the learning-grid
projection check separately. Fourier methods have direct spatial
Dysthe precedent \cite{fedele,lomei}; the specific gates here are proposed.}
\testdecision{Accept only cases passing the common reference budget and
spectral checks. A larger domain can alter nonlocal mean flow even when
edge amplitudes are tiny. An unresolved case triggers refinement or a
shorter qualified distance, not a claim of singularity. Full-period
wavetrains have a different boundary question from localized packets.}

\subsection{T2. Learn unseen initial fields at fixed parameters}
\testgoal{Determine whether the FNO learns the declared short-distance
operator rather than memorizing a small collection of trajectories.}
\testdesign{Train on the three-family group split at fixed epsilon and
length. Evaluate all held-out profiles without retraining. Compare
one-step and autonomous rollout errors to a coarser spectral solver
whose step and grid are tuned on validation for the same accuracy.
Include the exact linear propagator and independently resolved NLS
evolution as physical controls, while
recognizing that neither solves the full Dysthe problem.
Use nested training subsets to produce a learning curve \cite{fno}.}
\testdecision{Pass only if the prespecified held-out errors and failure
fraction meet the goal. A conservative first main-campaign target is
at least $90\%$ of test groups passing the primary gates through
$\Xi=4$, with uncertainty and all failures reported.
This establishes bounded-distribution interpolation only.
If a cheap spectral solve dominates accuracy and cost, report that
outcome before adding model complexity.}

\subsection{T3. Does the full physics loss actually help?}
\testgoal{Isolate the benefit, if any, of physics training for this
equation and input population.}
\testdesign{Compare the same initialized FNO backbone, data, hard initial
condition, linear skip, rollout schedule and validation protocol with
$\lambda_r=0$ and $\lambda_r>0$. Select loss weights using validation
only. Report both a matched-compute comparison and the extra cost to
reach a matched accuracy; physics updates cost more.
Evaluate on new collocation stations and refined full-profile residuals.
At a fixed architecture, repeat on two or three data budgets to test
data efficiency \cite{pino}.}
\testdecision{Physics is useful if it improves the intended held-out
metric or needs less reference data at acceptable total cost.
A lower training residual alone fails this decision.
A physics-only branch is optional after the hybrid works; it is not
necessary to launch a large architecture search. If benefits disappear
under independent residual refinement, investigate aliasing first.}

\subsection{T4. What can a PINN contribute?}
\testgoal{Separate single-problem optimization from genuine multi-input
physics-informed prediction.}
\testdesign{Choose six qualified validation cases, two per family.
For an ordinary Fourier PINN, compare a single interval, sequential
distance slabs, and slabs with causal residual weighting at comparable
optimization budgets \cite{causal,failures}. Carry predicted endpoints
forward; do not reset to exact interior snapshots.
If this reveals a useful training strategy, test a function-conditioned
Fourier/DeepONet model on unseen initial fields with no optimization
\cite{pideeponet}.}
\testdecision{Count per-case optimization, failed restarts and all slabs.
A PINN can be informative without beating the spectral solver.
The six-case exercise is a training diagnostic, not a broad statistical
comparison. A parameter-conditioned network only generalizes over its
declared packet parameters; an operator claim requires initial-field
conditioning and a disjoint field-level test set.}

\subsection{T5. How far does autonomous prediction remain useful?}
\testgoal{Measure a defensible accuracy horizon and whether rollout
training extends it.}
\testdesign{Train on reference intervals no longer than $\Xi=4$ and
repeat learned local maps on their own outputs to $\Xi=8,16$, then
conditionally 32. Compare one-step training, a short multi-step rollout
loss, and optionally stopped-gradient pushforward training
\cite{pushforward}. Use validation to select the rollout depth before
testing. Plot error and the qualified distance per held-out group.
Inspect spectral-error growth, not only norm drift \cite{stability}.}
\testdecision{Claim an extension only for the fixed test population and
tolerances. Never reset the rollout with a reference state.
For a locally Lipschitz exact map the conditional bound
$e_{n+1}\leq L_n e_n+\delta_n$ explains why small one-step errors
$\delta_n$ may accumulate. Here
$\delta_n=\norm{G_\theta(\widetilde u_n,\epsilon,\Delta)
-\mathcal S_\Delta^\epsilon(\widetilde u_n)}$ is the error
at a \emph{predicted} input; teacher-forced errors do not supply this
bound. This is not a proved uniform Dysthe estimate.
Do not enforce artificial contraction merely to keep output bounded.}

\subsection{T6. Preserve real modulational instability}
\testgoal{Determine whether a learned method reproduces physical
sideband amplification instead of damping it away.}
\testdesign{For amplitudes and modulation modes covering stable and
unstable regions, compare the full signed perturbation response with
$\exp(\xi M(k))$ from Appendix~\ref{app:mi}, a refined Dysthe trajectory,
and the NLS limit. Halve the perturbation until the early response
converges to the linear prediction while remaining above numerical
noise. Check both sideband amplitudes, their phases and growth rates.
The observable is motivated by Dysthe's instability calculation
\cite{dysthe1979}.}
\testdecision{Use a provisional $5\%$ gain/growth tolerance over the
verified linear interval, with absolute tolerances near zero growth.
Require the perturbation signal to exceed the reference error by a
declared margin. A filter that keeps a rollout smooth but suppresses
the correct growing mode fails. Linear amplification is neither
nonlinear blow-up nor wave breaking.}

\subsection{T7. Predict focusing and recurrence events}
\testgoal{Test the wave behavior that matters beyond average field error.}
\testdesign{For localized packets measure the first focus distance,
amplification, peak trajectory and asymmetry against refined Dysthe;
compare NLS to identify the size of the higher-order effect
\cite{lomei}. For periodic modulations measure carrier depletion,
sideband exchange and recurrence distance \cite{recurrence}.
Record raw complex recurrence defect and separately minimize over
translation and constant phase. If no event occurs in the qualified
reference interval, label it absent/censored rather than forcing a fit.}
\testdecision{Prespecify event definitions and the $5\%$ primary event
target. Similar peak oscillations do not establish recurrence of the
complex wave field or a stable breather. Track spectral mean and peak
separately and test the exact centroid identity
\eqref{eq:centroid-constraint}.
Two-packet collisions, with separated single-packet controls and a
larger-domain check, are a later family-transfer test; they are not
required to establish the first operator result.}

\subsection{T8. State exactly which generalization is being tested}
\testgoal{Find the boundary of the learned map's competence.}
\testdesign{Change one factor at a time: withheld packet family;
new epsilon values within the training interval; epsilon outside it;
finer grids on the same domain; a different domain length; longer
propagation. Each receives a separate score and qualification.
If testing measurement noise later, specify the observation and
reconstruction processes and compare to an appropriate estimator,
not to a solver supplied with inaccessible noiseless phase.}
\testdecision{An interpolation result does not imply extrapolation.
A finer output grid does not recover missing input frequencies.
A domain change alters physical wavenumbers and periodic mean flow.
Do not pool these tasks into one favorable average or declare
arbitrary-input performance.}

\subsection{T9. Is there an amortized computational benefit?}
\testgoal{Determine whether the learned solver saves work for the
intended repeated-query workload.}
\testdesign{Compare against the fastest numerical configuration that
passes the same field, derivative and event gates. Measure hardware,
precision, batch size, peak memory, output density, warm-up, inference,
data generation, validation, all training seeds and any test-time
optimization. Report accuracy--cost curves, not one speed ratio.
For a fixed workload a useful accounting identity is
\[
M_*=\frac{C_{\rm data}+C_{\rm train}}
          {C_{\rm numerical}-C_{\rm inference}},
\]
provided the denominator is positive; include other attributable
setup/selection costs in the numerator.}
\testdecision{The estimated break-even query count is meaningful only
at matched accuracy and for the specified hardware/batching regime.
If the denominator is nonpositive, there is no amortized inference
advantage under that comparison. A differentiable surrogate may still
be useful, but sensitivity accuracy then becomes a separate test.}
